% Samuel D. V. Stephens: one-page résumé for quantitative research and trading roles.
% Same type as the CV (Cormorant Garamond, from TeX Live by file name). Build with
%   python3 scripts/build-cv.py   (builds this and cv.tex; see the README)
\documentclass[11pt,letterpaper]{article}
\usepackage[margin=0.55in,top=0.45in,bottom=0.4in]{geometry}
\usepackage{fontspec}
\setmainfont{CormorantGaramond-Medium.ttf}[
    BoldFont=CormorantGaramond-Bold.ttf,
    ItalicFont=CormorantGaramond-MediumItalic.ttf,
    BoldItalicFont=CormorantGaramond-BoldItalic.ttf]
\newfontfamily\headingfont{CormorantGaramond-Bold.ttf}[LetterSpace=6.0]
\setmonofont{lmmono10-regular.otf}
\usepackage[USenglish]{babel}
\usepackage[strict,autostyle]{csquotes}
\usepackage[protrusion=true,expansion=false]{microtype}
\usepackage{enumitem}
\usepackage{xcolor}
\usepackage[hidelinks]{hyperref}
\hypersetup{pdftitle={Samuel D. V. Stephens: Résumé}, pdfauthor={Samuel D. V. Stephens}}

\pagestyle{empty}
\setlength{\parindent}{0pt}
\setlength{\parskip}{0pt}
\setlength{\emergencystretch}{1.5em}
\linespread{0.98}
\setlist[itemize]{leftmargin=1.1em, itemsep=0.5pt, topsep=1pt, parsep=0pt, label={\textbullet}}

% A section heading: letter-spaced small caps with a thin rule under it.
\newcommand{\sect}[1]{\vspace{4pt}{\headingfont\small\MakeUppercase{#1}}\par\vspace{1pt}\hrule height 0.4pt\vspace{3pt}}
% An entry: bold left text, date flush right; optional italic second line.
\newcommand{\entry}[2]{\textbf{#1}\hfill #2\par}
\newcommand{\role}[2]{\textit{#1}\hfill\textit{#2}\par}

\begin{document}

{\centering
{\fontsize{22}{24}\selectfont\bfseries Samuel D.\,V. Stephens}\par\vspace{3pt}
\href{mailto:samuelstephens2004@gmail.com}{samuelstephens2004@gmail.com}
\enspace·\enspace +1 214 843 5006
\enspace·\enspace \href{https://samuel-stephens.com}{samuel-stephens.com}
\enspace·\enspace \href{https://github.com/sdvstephens}{github.com/sdvstephens}\par}

\sect{Education}
\entry{Salem State University, Salem, MA}{2026--}
\role{B.S. in Mathematics and Music; GPA 4.00 (current)}{Expected 2029}
Graduate measure theory (directed study, Folland's \textit{Real Analysis}). Earlier coursework for credit at Wabash College (2025; Honors and Dean's Scholar) and Kenyon College (2023--24): complex analysis, abstract algebra, symbolic logic. Audited Harvard's Math 55a and 55b (honors abstract algebra; honors real and complex analysis), 2025--26.

\sect{Research Experience}
\entry{Stanford University School of Medicine}{2026--}
\role{Assistant Researcher, Department of Neurology}{Remote}
\begin{itemize}
  \item Extending CICADAS-II toward clinical use: personalized adaptive dosing thresholds from SOFA score and age, implemented against EEG and the MORGOTH large-scale ICU dataset.
\end{itemize}
\entry{Mass General Brigham}{2026--}
\role{Assistant Research Intern, Department of Neurology}{Boston, MA}
\begin{itemize}
  \item Real-data phase (CICADAS-III) on large ICU EEG and clinical datasets; first run of the framework on real patient records in October 2026.
\end{itemize}
\entry{Beth Israel Deaconess Medical Center \& Harvard Medical School}{2025--26}
\role{Assistant Researcher, Department of Neurology, under Dr.~M.~Brandon Westover}{Boston, MA}
\begin{itemize}
  \item Built a Monte Carlo simulator of 1,000 ICU patients: latent seizure-rate dynamics, sparse Poisson observations, one-compartment pharmacokinetics with bolus dosing, Hill-curve drug response, and a daily logistic mortality hazard.
  \item Estimated each patient's hidden seizure rate in real time with an extended Kalman filter on the log rate, and recovered causal effects of dosing strategies with the parametric g-formula where a naive comparison got the sign of the effect wrong.
  \item Derived a dosing rule that minimizes instantaneous marginal harm; untuned, it matched the best tuned fixed dose when one harm dominated and beat every tuned alternative when both harms were high.
  \item Stress-tested the estimator under misspecification: false detections, missing dose records, saturating harm, and unmeasured confounding (E-values). Co-first author of the resulting paper.
\end{itemize}

\sect{Publications and Manuscripts}
\begin{itemize}
  \item S.\,D.\,V. Stephens\textsuperscript{*}, T. Westover\textsuperscript{*}, S.\,F. Zafar, M.\,B. Westover. \enquote{CICADAS-II: Simulation-Based Causal Inference and Adaptive Dosing from Sparse Seizure Counts and Bolus Anti-Seizure Medication in the ICU.} Under major revision, \textit{Journal of Neural Engineering}. \textsuperscript{*}Co-first authors.
  \item S.\,D.\,V. Stephens. \textit{Studies in Algebra and Group Theory}. 519-page textbook for Harvard's Math 55a; Professor Joe Harris is offering it as supplementary material to his 2026--27 students.
  \item S.\,D.\,V. Stephens. \enquote{Monotone Arithmetic Progressions in the Gray Code and in Negabinary Enumerations of the Integers} (Erdős Problems 195 and 196). In preparation.
\end{itemize}

\sect{Projects}
\begin{itemize}
  \item \textbf{cicadas2} (MATLAB; Python with NumPy, SciPy, pytest). Primary developer of the paper's reference implementation; Python port held to the MATLAB original by cross-implementation equivalence tests, and a harness that regenerates every reported number. Private until publication.
  \item \textbf{rmpp-stabilizer} (C++). \texttt{LD\_PRELOAD} library that intercepts tablet pen input before the system UI and smooths it in real time (moving-average, Gaussian, string-pull and One Euro filters); Docker cross-compilation toolchain for ARM64.
  \item \textbf{acad-cli} (Python). Command-line coursework manager published on PyPI (\texttt{pipx install acad-cli}): \LaTeX{} and Org document generation, Taskwarrior integration, plain-JSON records under Git.
\end{itemize}

\sect{Skills}
\textbf{Languages:} Python (NumPy, SciPy, pytest), MATLAB, C++, R, Mathematica, \LaTeX.\enspace
\textbf{Tools:} Git, Linux and the Unix shell, Docker, Neovim.\enspace
\textbf{Methods:} state-space models and Kalman filtering, Monte Carlo simulation, causal inference, survival and hazard models, PK/PD modeling.

\sect{Other}
Private mathematics tutor (precalculus, AP Calculus BC, physics), 2025--.\enspace KEEP-STEM Scholar, Kenyon College program supported by the National Science Foundation (2023).

\end{document}
